\chapter{Beweisbilanz und Anschluss an die bisherigen Bände}

\section{Zielaussagen und noch offene Beweisketten}

\noindent\textbf{Vorbehalt fuer die folgende Uebersicht.}
Die angegebenen Zielaussagen duerfen noch nicht als lueckenlos im
vereinbarten Format bewiesen gelten. Die Existenz und Eindeutigkeit
der Erfuellungsrelation sind jetzt in
\FormulaRefAuto{B48SatisfactionExists}[thm] mit konkreter Baumcodierung,
Strukturrekursion und einem getrennten Eindeutigkeitsbeweis ausgearbeitet.
Das Koinzidenzlemma \FormulaRefAuto{B48Coincidence}[thm], das
Substitutionslemma \FormulaRefAuto{B48SubstitutionLemma}[thm] und
die Korrektheit elementarer Herleitungen
\FormulaRefAuto{B48Soundness}[thm] sind nun ausgearbeitet.
Der Korrektheitsbeweis behandelt alle 21 elementaren Regelfaelle
einschliesslich Annahmen, Entladungen und Eigenvariablenbedingungen.
Die Henkin-Konstruktion und die dafuer benoetigten syntaktischen
Beweisumformungen in Kapitel~2 bleiben dagegen offen.
Ausgearbeitet sind ausserdem einzelne semantische Schlussformen
mit expliziten Regel- und Zeilenverweisen. Die anschliessenden
Unabhaengigkeitsargumente erben die offenen Voraussetzungen.

\begin{longtable}{@{}p{0.42\linewidth}p{0.53\linewidth}@{}}
\toprule
\textbf{Aussage} & \textbf{Status und Beweisstelle}\\
\midrule
\endhead
\(A\approx B\Rightarrow\mathcal P(A)\approx\mathcal P(B)\)
 & Beweisskizze mit Bild- und Urbildabbildung:
 \FormulaRefAuto{B48PowerTransport}[thm].\\[5pt]
\(|A|<|\mathcal P(A)|\)
 & Beweisskizze mit der Diagonalmenge:
 \FormulaRefAuto{B48Cantor}[thm].\\[5pt]
Kürzung bei einer endlichen Ausgangsmenge
 & Tabellenbeweis noch auszuarbeiten:
 \FormulaRefAuto{B48FiniteCancellation}[thm].\\[5pt]
Globale Potenzmengenkürzung PK
 & Ziel: unter \(\operatorname{Con}(\mathsf{ZFC})\) weder in ZFC
 beweisbar noch widerlegbar; die Voraussetzungen sind noch offen:
 \FormulaRefAuto{B48PCIndependence}[thm].\\[5pt]
\(|\mathbb R|=2^{\aleph_0}\)
 & Beweisskizze mit zwei Injektionen und
 Cantor--Schröder--Bernstein:
 \FormulaRefAuto{B48RealContinuum}[thm].\\[5pt]
CH als Fehlen einer Zwischenmächtigkeit
 & Beweisskizze zur Aequivalenz in ZFC:
 \FormulaRefAuto{B48CHIntermediate}[thm].\\[5pt]
Kontinuumshypothese CH
 & Ziel: unter \(\operatorname{Con}(\mathsf{ZFC})\) weder in ZFC
 beweisbar noch widerlegbar; die Voraussetzungen sind noch offen:
 \FormulaRefAuto{B48CHIndependence}[thm].\\
\bottomrule
\end{longtable}

\section{Axiomatische Abhängigkeiten}

Die erste Ebene bleibt die Mengenlehre aus Band~3 mit den bereits
entwickelten Funktionen, Bijektionen, Wohlordnungen und Zahlbereichen.
Darauf baut dieser Band die folgende Kette auf:
\begin{longtable}{@{}p{0.44\linewidth}p{0.51\linewidth}@{}}
\toprule
\textbf{Grundlage} & \textbf{Darauf aufbauendes Beweisziel}\\
\midrule\endhead
Formelcodes, Erfuellungsrelation, Koinzidenz und Substitution
 & Korrektheit (ausgearbeitet), Henkin-Vollstaendigkeit (noch offen)\\[0.5ex]
Ordinale und Rekursion & Ränge, Skolemhüllen und Kollaps\\[0.5ex]
\(L\)-Hierarchie, Reflexion und Kondensation
 & Relative Konsistenz von GCH\\[0.5ex]
Namen, Forcingrelation und generischer Quotient
 & Erhaltung von ZFC\\[0.5ex]
ccc und Zählung schöner Namen
 & \(2^{\aleph_0}=2^{\aleph_1}=\aleph_2\)\\
\bottomrule
\end{longtable}
Aus den beiden Modellmöglichkeiten folgt mittels Korrektheit und
Deduktionskriterium die jeweilige Unabhängigkeit.
Keiner der beiden Modellwege wird als zusätzliches Mengenaxiom
vorausgesetzt. Auch ein abzählbares transitives Modell ist keine
verdeckte Voraussetzung.

Der Goedel'sche Unvollstaendigkeitssatz muss nicht vorangestellt
werden. Der hier benoetigte Vollstaendigkeitssatz ist in Kapitel~2
noch nicht vollstaendig ausgearbeitet. Die ausdrueckliche
Konsistenzannahme bleibt unveraendert bestehen.

\section{Reichweite und Darstellungsniveau}

\paragraph{Umstellungsstand der Begr\"undungsspalten.}
Die neuen Modellhilfss\"atze, das Unabh\"angigkeitskriterium und die
abschlie\ss enden Beweise der Nicht-Beweisbarkeit, Nicht-Widerlegbarkeit
und Unabh\"angigkeit sind mit konkreten Regel-, Satz- und Zeilenverweisen
aufgeschrieben. Diese Umstellung ist noch nicht f\"ur alle Beweise des
Bandes abgeschlossen. Insbesondere sind die weiterhin frei formulierten
Begr\"undungen in den Grundlagen-, Konstruktibilit\"ats- und Forcingbeweisen
noch einzeln auf zuvor bewiesene Hilfss\"atze zur\"uckzuf\"uhren. Die
folgende Beweisbilanz ist daher keine Behauptung einer bereits l\"uckenlos
referenzierten Ableitung s\"amtlicher Voraussetzungen.


Der Band wiederholt nicht die gesamte axiomatische Mengenlehre.
Er ergänzt genau die Modell- und Hierarchiemethoden, welche
die Zielkapitel benötigen. Die Beweistabellen enthalten
mathematische Ableitungen mit definierten Begriffen und
Verweisen auf Hilfssätze. Ein vorhandener Verweis allein besagt
jedoch nicht, dass der zitierte Hilfssatz schon im verlangten Format
bewiesen ist; insbesondere bleiben die oben genannten Grundlagen offen.
Eine Übersetzung aller Tabellen in einen maschinell überprüften
Beweiskalkül wäre ein eigenes Formalisierungsprojekt; die
vorliegende Darstellung beansprucht dies nicht.

Weitergehende Themen wie die allgemeine Easton-Theorie,
iteriertes Forcing, große Kardinalzahlen oder die
Unvollständigkeitssätze sind für die hier erreichten Resultate
nicht erforderlich und werden deshalb nicht als weitere
Pflichtbände eingeführt.

\section{Literatur zur Einordnung der Konstruktionen}

Die folgenden frei zugänglichen Fachtexte dienen als
weiterführende Belege und zur Vertiefung der klassischen Methoden.
Die Definitionen und Beweisschritte, die in den Kapiteln~4 und~5
benutzt werden, sind im vorliegenden Band selbst angegeben.

\begin{itemize}
 \item Andrew Marks:
 \emph{Set Theory}, Vorlesungsnotizen, Fassung vom 31.~Januar 2026.
 Kapitel über konstruktible Mengen, Kondensation, GCH und Forcing.
 \href{https://math.berkeley.edu/~marks/notes/set_theory_notes_4.pdf}{Vorlesungsnotizen (Berkeley).}
 \item Chris Lambie-Hanson:
 \emph{The Continuum Hypothesis, the Axiom of Choice, and
 Lebesgue Measurability}.
 \href{https://www.math.cmu.edu/~clambieh/ch.pdf}{Vortrag (Carnegie Mellon University).}
 \item Regula Krapf:
 \emph{Forcing: How to prove unprovability -- The Continuum Hypothesis}.
 Darstellung des Cohen-Forcings und der Rolle der
 abzählbaren Kettenbedingung.
 \href{https://www.math.uni-bonn.de/people/krapf/ForcingCHBielefeld.pdf}{Vortrag über Forcing und CH (Universität Bonn).}
\end{itemize}
